\section{Conclusion and Future Work}

We presented library learning modulo theory (LLMT),
 a technique for learning abstractions from a corpus of programs
 modulo a user-provided equational theory.
We implemented LLMT in \tool.
Our evaluation showed that \tool
 achieves better compression orders of magnitude faster
 than the state of the art.
On a larger benchmark suite of \cogsci programs,
 \tool learns sensible functions that
 compress a dataset that was---until now---too large for library learning techniques.

LLMT and \tool present many avenues for future work.
First,
 our evaluation showed that equational theories are important
 for achieving high compression,
 but these must be provided by domain experts.
Recent work in automated theory synthesis like
 Ruler~\cite{ruler} or \tname{TheSy}~\cite{thesy}
 could aid the user in this task.
Second,
 LLMT uses e-graph anti-unification to generate
 promising abstraction candidates, 
 but this approach is incomplete and misses some patterns
 that could achieve better compression.
% Our prototype indicated that 
%  a complete approach based on traversing the
%  pattern generalization lattice 
%  was prohibitively slow;
%  future work could identify a subset of patterns
%  that does not sacrifice optimality 
%  but can still be enumerated efficiently. 
An exciting direction for future work is to combine LLMT 
with more efficient top-down search from \stitch~\cite{stitch}.
%
This is challenging because \stitch crucially relies on the ability 
to quickly compute an upper bound on the compression of a given pattern 
by summing up the local compression at each of its matches in the corpus.
%
This upper bound does not straightforwardly extend to e-graphs
because in an e-graph different matches of a pattern may come from different syntactic variants of the corpus,
and one needs to trade-off the compression from abstractions
against the size difference between different syntactic variants.